\section{Discussion}
\label{sec:discussion}

\subsection{Key Findings and Their Implications}

\textbf{SE and min\_logprob measure genuinely distinct uncertainty dimensions.} The Pearson $|r|=0.026$ is not merely a marginal pass of the 0.70 threshold --- it indicates near-orthogonality by any reasonable standard. This gap from the prior first-token measurement ($r=0.54$--$0.76$) is mechanistically explained: min\_logprob captures late-sequence token uncertainty while first-token confidence captures only the initial prediction. The operational separation between SE (stochastic sampling, semantic aggregation) and min\_logprob (greedy decode, token-level minimum) manifests as near-zero empirical correlation.

This finding establishes a necessary precondition for ensemble benefit: it rules out the scenario where combining SE and min\_logprob provides no benefit because one is a redundant reparameterization of the other. However, necessity is not sufficiency --- and the partial $R^2$ null result ($0.0005$ at $N=2500$) shows that independence does not automatically yield linear ensemble gain. The signals are empirically orthogonal yet carry redundant correctness-predictive information in the linear logistic regression setting.

\textbf{Cross-model judge design suppresses circularity effectively.} The $\rho(\text{SE}, \text{judge}) = -0.026$ result validates the evaluation protocol: using Qwen-2.5-7B to evaluate Llama-3.1-8B outputs prevents shared model-specific NLI biases from creating circular agreement. This design choice, recommended by \citet{santilli2025revisiting}, is empirically validated in our setting.

\textbf{The partial $R^2$ null result: power characterization and gate disclosure.} The partial $R^2=0.0005$ at $N=2500$ (LRT $p=0.442$) was classified by the pre-registered pipeline gate as \texttt{EXPLORE\_N10} (gate\_pass: false) --- the gate did not pass, and its EXPLORE\_N10 recommendation motivated the h-e1-v2 follow-up to extend to $N=10$ stochastic samples. We disclose this gate outcome transparently.

At the same time, we characterize the null finding as adequately powered for the pre-registered effect size. $N=2500$ substantially exceeds the sample size required for $>0.80$ power to detect the pre-registered effect size of partial $R^2\geq0.02$ (estimated power $>0.99$). The gate's conservative EXPLORE\_N10 flag reflects the pipeline's joint marginal-uncertainty criterion --- it fires when either the Pearson $r$ or the partial $R^2$ falls in a marginal zone --- not a judgment that the sample size is insufficient to detect the pre-registered threshold. We therefore characterize the partial $R^2=0.0005$ finding as a powered null for the pre-registered effect size: not an absence of evidence due to underpowering, but evidence of absence of a linear conditional effect at the pre-registered threshold on this benchmark. The dissociation between near-zero correlation ($|r|=0.026$) and near-zero conditional contribution (partial $R^2=0.0005$) suggests that, while the signals are computed differently, they capture redundant correctness-predictive information in the linear setting.

\subsection{Limitations}

\textbf{Ensemble AUROC not measured; linear null bounds linear gain.} Ensemble AUROC is not evaluated in this work. The partial $R^2$ null at $N=2500$ (gate$=$EXPLORE\_N10) suggests linear logistic regression ensembles will not achieve the $\Delta$AUROC target; $\Delta$AUROC $\geq 0.025$ on TriviaQA and cross-model transfer to Qwen-2.5-7B on TruthfulQA are neither confirmed nor refuted and are the primary targets for h-e1-v2, which will also extend to $N=10$ stochastic samples per prompt as the gate recommends.

\textbf{Single seed and no bootstrapped confidence intervals.} $N=2500$ with a fixed seed (42) provides no estimate of result variance. The Pearson $|r|=0.026$ and Spearman $\rho=-0.026$ are point estimates; bootstrapped 95\% CIs at $N=2500$ are computationally cheap and are planned for the h-e1-v2 analysis alongside nonlinear ensemble evaluation.

\textbf{Scope: short-answer factual QA with 7--8B open-weight models.} Results are validated on TriviaQA dev ($\leq$50 token answers) with Llama-3.1-8B-Instruct. Generalization to long-form generation is explicitly out of scope --- the SNNE paper \citep{nguyen2025beyond} documents SE's weakness for multi-sentence answers.

\subsection{Broader Impact}

This work contributes to the safety and reliability of LLM deployments by establishing the empirical independence of two lightweight UQ signals --- one requiring only greedy decode access (min\_logprob) and one requiring only generation API access ($\text{SE}_{N5}$). The replicable evaluation methodology (cross-model LM-judge, circularity diagnostic, mechanism reality checks) provides a template for future UQ signal characterization studies. We do not foresee significant negative impacts: better hallucination detection reduces risks in high-stakes settings and does not enable new attack vectors. The pipeline uses publicly available open-weight models, ensuring reproducibility without proprietary API dependence.
